\textbf{Phần 1. Trắc nghiệm lựa chọn}
\setcounterpageref{ex}{0}

\begin{ex} %[3D1V06]
Tính giá trị của biểu thức: $A = 4 \times 8 - 36 : 4 + 12$.
\choice
{$33$}
{\True $35$}
{$37$}
{$31$}
\loigiai{
Thực hiện nhân và chia trước, cộng trừ sau:\\
$$A = 32 - 9 + 12 = 23 + 12 = 35$$
\textbf{Đáp án của câu hỏi là B}
}
\end{ex}

\begin{ex} %[3D1V06]
Một đàn bò có $7$ con và một đàn lợn có $5$ con. Hỏi cả đàn bò và đàn lợn đó có tất cả bao nhiêu cái chân?
\choice
{\True $48$ cái chân}
{$44$ cái chân}
{$52$ cái chân}
{$40$ cái chân}
\loigiai{
Mỗi con bò và mỗi con lợn đều có $4$ cái chân.\\
Tổng số con bò và lợn là:
$$7 + 5 = 12\text{ (con)}$$
Cả hai đàn có tất cả số cái chân là:
$$12 \times 4 = 48\text{ (cái chân)}$$
\textbf{Đáp án của câu hỏi là A}
}
\end{ex}

\begin{ex} %[3D1V06]
Tìm giá trị của $x$, biết: $x : 4 + 23 = 32$.
\choice
{$32$}
{$40$}
{\True $36$}
{$28$}
\loigiai{
Xem $x : 4$ là số hạng chưa biết:\\
$$x : 4 = 32 - 23$$
$$x : 4 = 9$$
$$x = 9 \times 4 = 36$$
\textbf{Đáp án của câu hỏi là C}
}
\end{ex}

\begin{ex} %[3D1V06]
Một chiếc khăn tay hình vuông có cạnh dài $8\text{ cm}$. Người ta viền ren xung quanh chiếc khăn tay đó. Độ dài đoạn ren cần dùng là
\choice
{$24\text{ cm}$}
{$28\text{ cm}$}
{$36\text{ cm}$}
{\True $32\text{ cm}$}
\loigiai{
Độ dài đoạn ren viền xung quanh khăn tay chính là chu vi hình vuông:\\
$$8 \times 4 = 32\text{ (cm)}$$
\textbf{Đáp án của câu hỏi là D}
}
\end{ex}

\begin{ex} %[3D1V06]
Điền dấu thích hợp vào chỗ chấm: $4 \times 6 - 8 \dots 32 : 4 \times 2$.
\choice
{$>$}
{\True $=$}
{$<$}
{Không so sánh được}
\loigiai{
Tính giá trị hai vế:\\
- Vế trái: $4 \times 6 - 8 = 24 - 8 = 16$\\
- Vế phải: $32 : 4 \times 2 = 8 \times 2 = 16$\\
Vì $16 = 16$ nên ta điền dấu $=$.\\
\textbf{Đáp án của câu hỏi là B}
}
\end{ex}

\begin{ex} %[3D1V06]
Trong một phép chia cho $4$, có thương bằng $6$ và số dư là số dư lớn nhất có thể có. Số bị chia của phép chia đó là
\choice
{\True $27$}
{$26$}
{$28$}
{$25$}
\loigiai{
Vì số chia là 4 nên số dư lớn nhất có thể có là $3$.\\
Số bị chia là:
$$4 \times 6 + 3 = 24 + 3 = 27$$
\textbf{Đáp án của câu hỏi là A}
}
\end{ex}

\begin{ex} %[3D1V06]
Khối lớp 3 có $36$ bạn học sinh tham gia đồng diễn thể dục. Thầy giáo chia đều các bạn thành $4$ hàng, sau đó mỗi hàng lại chia thành $3$ nhóm nhỏ. Hỏi mỗi nhóm nhỏ có bao nhiêu bạn học sinh?
\choice
{$4$ bạn}
{$2$ bạn}
{\True $3$ bạn}
{$5$ bạn}
\loigiai{
Mỗi hàng có số bạn học sinh là:
$$36 : 4 = 9\text{ (bạn)}$$
Mỗi nhóm nhỏ có số bạn học sinh là:
$$9 : 3 = 3\text{ (bạn)}$$
\textbf{Đáp án của câu hỏi là C}
}
\end{ex}

\begin{ex} %[3D1V06]
Tính nhanh giá trị của biểu thức: $4 \times 1 + 4 \times 2 + 4 \times 3 + 4 \times 4$.
\choice
{$36$}
{$44$}
{$48$}
{\True $40$}
\loigiai{
Đặt thừa số chung:\\
$$4 \times (1 + 2 + 3 + 4) = 4 \times 10 = 40$$
\textbf{Đáp án của câu hỏi là D}
}
\end{ex}

\begin{ex} %[3D1V06]
Tìm giá trị của $x$, biết: $(x - 7) \times 4 = 28$.
\choice
{$12$}
{\True $14$}
{$16$}
{$15$}
\loigiai{
Xem $(x - 7)$ là thừa số chưa biết:\\
$$x - 7 = 28 : 4$$
$$x - 7 = 7$$
$$x = 7 + 7 = 14$$
\textbf{Đáp án của câu hỏi là B}
}
\end{ex}

\begin{ex} %[3D1V06]
Bác Ba nuôi $6$ con thỏ. Số con gà bác Ba nuôi gấp $4$ lần số con thỏ. Hỏi bác Ba nuôi số con gà nhiều hơn số con thỏ bao nhiêu con?
\choice
{\True $18$ con}
{$24$ con}
{$16$ con}
{$20$ con}
\loigiai{
Số con gà bác Ba nuôi là:
$$6 \times 4 = 24\text{ (con)}$$
Số con gà nhiều hơn số con thỏ là:
$$24 - 6 = 18\text{ (con)}$$
\textbf{Đáp án của câu hỏi là A}
}
\end{ex}

\begin{ex} %[3D1V06]
Cho dãy số cách đều: $4; 8; 12; 16; \dots$ Số hạng thứ 8 của dãy số trên là
\choice
{$28$}
{$36$}
{\True $32$}
{$40$}
\loigiai{
Dãy số gồm các số chia hết cho 4. Số hạng thứ 8 là:
$$4 \times 8 = 32$$
\textbf{Đáp án của câu hỏi là C}
}
\end{ex}

\begin{ex} %[3D1V06]
Trong một bãi đỗ xe có $8$ chiếc xe ô tô con (mỗi xe có $4$ bánh) và $5$ chiếc xe máy (mỗi xe có $2$ bánh). Tổng số bánh xe của cả hai loại xe đó là
\choice
{$38$ bánh xe}
{$40$ bánh xe}
{$44$ bánh xe}
{\True $42$ bánh xe}
\loigiai{
Số bánh xe của 8 chiếc ô tô con là:
$$8 \times 4 = 32\text{ (bánh xe)}$$
Số bánh xe của 5 chiếc xe máy là:
$$5 \times 2 = 10\text{ (bánh xe)}$$
Tổng số bánh xe là:
$$32 + 10 = 42\text{ (bánh xe)}$$
\textbf{Đáp án của câu hỏi là D}
}
\end{ex}

\textbf{Phần 2. Trắc nghiệm đúng sai}
\setcounterpageref{ex}{0}

\begin{ex} %[3D1V06]
Xét tính đúng sai của các khẳng định sau về phép nhân và phép chia với $4$:
\choiceTF
{\True Tích của số $4$ với một số tự nhiên bất kì luôn luôn là một số chẵn.}
{Trong một phép chia cho $4$, số dư có thể bằng $4$.}
{\True Một số tự nhiên chia hết cho $4$ thì chắc chắn chia hết cho $2$.}
{\True Giá trị của biểu thức $4 \times 9$ lớn hơn giá trị của biểu thức $5 \times 7$.}
\loigiai{
\begin{itemize}
    \item Ý a) Đúng: Vì $4 = 2 \times 2$ nên $4 \times n = 2 \times (2 \times n)$ luôn là số chẵn.
    \item Ý b) Sai: Số dư luôn phải nhỏ hơn số chia (số dư có thể là 0, 1, 2, 3; không thể bằng 4).
    \item Ý c) Đúng: Vì 4 chia hết cho 2 nên số chia hết cho 4 cũng chia hết cho 2.
    \item Ý d) Đúng: $4 \times 9 = 36$ và $5 \times 7 = 35$. Vì $36 > 35$ nên nhận định đúng.
\end{itemize}
\begin{itemize}
    \item a) \textbf{Đúng}.
    \item b) \textbf{Sai}.
    \item c) \textbf{Đúng}.
    \item d) \textbf{Đúng}.
\end{itemize}
}
\end{ex}

\begin{ex} %[3D1V06]
Một cuộn dây thừng dài $36\text{ m}$ được cắt thành $4$ đoạn dài bằng nhau:
\choiceTF
{\True Mỗi đoạn dây dài $9\text{ m}$.}
{Để cắt cuộn dây thành $4$ đoạn, người ta cần thực hiện đúng $4$ nhát cắt.}
{\True Tổng độ dài của $3$ đoạn dây là $27\text{ m}$.}
{\True Dùng $1$ đoạn dây đó để viền quanh một khung tranh hình vuông có cạnh dài $2\text{ m}$ thì còn thừa đúng $1\text{ m}$ dây.}
\loigiai{
\begin{itemize}
    \item Ý a) Đúng: Chiều dài mỗi đoạn là: $36 : 4 = 9\text{ (m)}$.
    \item Ý b) Sai: Để cắt thành 4 đoạn chỉ cần $4 - 1 = 3$ nhát cắt.
    \item Ý c) Đúng: 3 đoạn dây dài: $3 \times 9 = 27\text{ (m)}$.
    \item Ý d) Đúng: Chu vi khung tranh là: $2 \times 4 = 8\text{ (m)}$. Đoạn dây 9 m còn thừa: $9 - 8 = 1\text{ (m)}$.
\end{itemize}
\begin{itemize}
    \item a) \textbf{Đúng}.
    \item b) \textbf{Sai}.
    \item c) \textbf{Đúng}.
    \item d) \textbf{Đúng}.
\end{itemize}
}
\end{ex}

\textbf{Phần 3. Trả lời ngắn}
\setcounterpageref{ex}{0}

\begin{ex} %[3D1V06]
Tính giá trị của biểu thức: $4 \times 7 - 20 : 4$.
\shortans{23}
\loigiai{
\textbf{Đáp số:} $23$\\
Thực hiện nhân chia trước:\\
$$4 \times 7 - 20 : 4 = 28 - 5 = 23$$
}
\end{ex}

\begin{ex} %[3D1V06]
Một trang trại nuôi $9$ con thỏ. Hỏi $9$ con thỏ đó có tất cả bao nhiêu cái chân?
\shortans{36}
\loigiai{
\textbf{Đáp số:} $36$\\
Mỗi con thỏ có 4 cái chân. Số cái chân của 9 con thỏ là:
$$9 \times 4 = 36\text{ (cái chân)}$$
}
\end{ex}

\begin{ex} %[3D1V06]
Tìm số bị chia trong một phép chia cho $4$, biết thương bằng $7$ và số dư bằng $2$.
\shortans{30}
\loigiai{
\textbf{Đáp số:} $30$\\
Số bị chia là:
$$4 \times 7 + 2 = 28 + 2 = 30$$
}
\end{ex}

\begin{ex} %[3D1V06]
Lớp 3C có $28$ học sinh được chia đều thành $4$ tổ. Hỏi $3$ tổ của lớp 3C có tất cả bao nhiêu học sinh?
\shortans{21}
\loigiai{
\textbf{Đáp số:} $21$\\
Mỗi tổ có số học sinh là:
$$28 : 4 = 7\text{ (học sinh)}$$
3 tổ có tất cả số học sinh là:
$$7 \times 3 = 21\text{ (học sinh)}$$
}
\end{ex}

\begin{ex} %[3D1V06]
Tìm giá trị của $x$, biết: $x : 4 = 15 - 8$.
\shortans{28}
\loigiai{
\textbf{Đáp số:} $28$\\
Tính vế phải:\\
$$x : 4 = 7$$
$$x = 7 \times 4 = 28$$
}
\end{ex}

\begin{ex} %[3D1V06]
Tính nhanh giá trị của biểu thức: $4 \times 2 + 4 \times 3 + 4 \times 5$.
\shortans{40}
\loigiai{
\textbf{Đáp số:} $40$\\
Đặt thừa số chung:\\
$$4 \times (2 + 3 + 5) = 4 \times 10 = 40$$
}
\end{ex}
